\chapter{Introducción}
\label{cap:introduccion}

Los criptoactivos se negocian en intercambios centralizados (CEX) y descentralizados (DEX), incluidos los creadores automáticos de mercado (AMM), cuyos precios cotizados dependen de las reservas de la piscina. El arbitraje puede acercar el precio de un AMM a una referencia externa, pero la magnitud y velocidad del ajuste dependen de la regla de intercambio, los costos, las reservas disponibles, el capital y el orden de las transacciones. Estos mecanismos se estudian en la literatura sobre pérdida frente al rebalanceo (LVR), arbitraje y valor extraíble por mineros o máximo (MEV) \cite{milionis2022lvr,daian2020flash,zhou2021hft}. Esta tesis aísla un entorno mucho más acotado: una piscina de producto constante sin comisiones, una referencia CEX exógena y una regla de arbitraje limitado.

La primera contribución es analítica. Un intercambio de $c$ veces la reserva de entrada cambia el precio logarítmico en $2\log(1+c)$ hacia la referencia, hasta alcanzar la banda de no arbitraje. Esto permite calcular los ticks ideales hasta la banda sin simulación. La pregunta computacional es qué cambia cuando actúan el ruido, la tendencia y un disparador de retiro único de liquidez en un calendario fijo y durante un horizonte finito de eventos.

El modelo basado en agentes es deliberadamente pequeño y reproducible. Representa una piscina de producto constante, una referencia exógena y reglas estilizadas de ruido, tendencia, arbitraje y proveedor de liquidez. La regla de retiro se activa en todas las ejecuciones de la condición adaptativa ensayada; por eso se trata como una intervención fija y no como evidencia de adaptación heterogénea de proveedores. En el diseño principal varía la profundidad y se mantienen fijos los tamaños absolutos de órdenes; un análisis separado los escala con las reservas. Por tanto, el experimento muestra cómo las conclusiones dependen de la razón entre el tamaño de orden y las reservas y de la regla de pánico, no cómo la profundidad determina por sí sola la respuesta de un mercado real.

El análisis de modelos reducidos compara persistencia, un AR(1) escalar, un mapa afín de dos estados y un mapa sensible a la banda en trayectorias del simulador. El mapa de dos estados no supera al autorregresivo escalar. La comparación del mapa sensible a la banda ofrece una mejora pequeña, concentrada en una condición sintética, y recibe como entrada la trayectoria futura realizada del CEX; es un diagnóstico de reproducción, no un pronóstico en tiempo real. Los resultados no sustentan afirmaciones de clasificación de regímenes ni de estabilidad.

El hallazgo central con valor para futuras investigaciones es que la primera recuperación puede ser temporal: el 31.2\,\% de las trayectorias simuladas que volvió a la banda de no arbitraje salió de ella después; la proporción aumenta a 98.8\,\% con la menor profundidad ensayada. Esto motiva distinguir el primer retorno del seguimiento sostenido en estudios futuros de reproducción de eventos. La tesis aporta un experimento reproducible centrado en mecanismos y una separación clara entre consecuencias algebraicas y resultados de simulación dependientes de reglas; no establece efectos en mercados CEX/DEX en operación.

El documento se organiza así. El capítulo~\ref{cap:problema} delimita la pregunta de investigación y revisa antecedentes; el capítulo~\ref{cap:justificacion} explica el aporte científico y metodológico. Los capítulos~\ref{cap:hipotesis} y~\ref{cap:objetivos} presentan la proposición analítica, las preguntas de evaluación y los objetivos. El capítulo~\ref{cap:metodologia} especifica el modelo y el diseño de simulación; el capítulo~\ref{cap:resultados} informa los hallazgos analíticos y computacionales, y el capítulo~\ref{cap:discusion} discute su alcance. El capítulo~\ref{cap:futuro} expone el trabajo necesario para reproducir eventos empíricos y ampliar el modelo; la bibliografía y los anexos técnicos cierran el documento.
